\section{Observation rules for complete paths}
\label{sec:observations}

Three observation rules act differently on a complete ordered path.
A downstream angle retains a whole contact and a following free prefix;
first hitting the contact angle cuts at admission and is discontinuous
at that path; a fixed local patch crossed strictly by the contact selects
any prescribed interior prefix.

% SCHEMATIC projection of one reference contact to (x=-p/a,Psi).
% Other state coordinates are suppressed. No finite-epsilon contact is drawn.
\begin{figure}[!htbp]
\centering
\begin{tikzpicture}[
 x=1cm,y=1cm,font=\small,
 every node/.style={inner sep=1.3pt},
 arrow/.style={-{Latex[length=1.5mm,width=1mm]},line width=.5pt},
 path/.style={draw=black,line width=.8pt},
 section/.style={draw=black!50,densely dashed,line width=.55pt}
]
 \foreach \shift in {0,5.0,10.0}{
  \begin{scope}[xshift=\shift cm]
   \draw[arrow] (.15,.35) -- (4.35,.35) node[right] {$x$};
   \draw[arrow] (.15,.35) -- (.15,3.80) node[above] {$\Psi$};
   \draw[path] (.45,1.40) .. controls (.95,1.18) and (1.25,1.07) .. (1.75,1.02);
   \draw[path] (1.75,1.02) -- (1.75,2.95);
   \draw[arrow] (1.75,1.55) -- (1.75,2.35);
   \draw[path,arrow] (1.75,2.95) .. controls (2.35,2.90) and (2.90,2.78) .. (3.75,2.70);
   \draw[fill=white,line width=.7pt] (1.75,1.02) circle[radius=2pt];
   \draw[fill=white,line width=.7pt] (1.75,2.95) circle[radius=2pt];
   \node[below=3pt] at (1.75,.35) {$s$};
  \end{scope}
 }
 \node[font=\small\bfseries] at (2.25,4.55) {(a) Downstream section};
 \node[font=\small\bfseries] at (7.25,4.55) {(b) Fixed-angle first stop};
 \node[font=\small\bfseries] at (12.25,4.55) {(c) Local crossed patch};
 \draw[section] (3.35,.70) -- (3.35,3.70);
 \node[below=3pt] at (3.35,.35) {$\sigma$};
 \fill (3.35,2.74263) circle[radius=2pt];
 \node[align=center] at (2.25,-.35)
   {common $\varepsilon_m(s)$\\complete contact retained};
 \draw[section] (6.75,.70) -- (6.75,3.70);
 \node[anchor=west] at (7.02,1.02) {$\mathcal T_s\Gamma$};
 \fill (6.67,2.87) rectangle (6.83,3.03);
 \node[anchor=west,align=left] at (7.02,3.35)
   {$\lim\mathcal T_s f_{\varepsilon_\ell}^{\sigma}$};
 \node[align=center] at (7.25,-.35)
   {a separate diagonal sequence\\first-hit limits differ};
 \draw[draw=black!40,densely dotted,line width=.5pt]
   (11.75,2.08) circle[radius=.50cm];
 \draw[line width=.8pt] (11.25,2.08) -- (12.25,2.08);
 \fill (11.75,2.08) circle[radius=2pt];
 \node[anchor=west] at (12.48,2.08) {$\Sigma$};
 \node[anchor=west] at (12.02,2.53) {$X_*$};
 \node[anchor=east] at (11.43,1.70) {$\alpha_0$};
 \node[align=center] at (12.25,-.35)
   {the same common $\varepsilon_m(s)$\\prescribed interior prefix};
\end{tikzpicture}
\caption{Three observations of one limiting contact, schematically
projected to $x=-p/a$ and $\Psi=v-c_\Sigma u_0$
(Corollary~\ref{cor:common-local-contact}), with $s<\sigma<F_\lambda(s)$.
Open circles mark admission and release; arrows indicate curve order.
(a) A downstream section retains the whole contact along the common
sequence $\varepsilon_m(s)$.
(b) The limit path $\Gamma$ first reaches $x=s$ at admission, while
actual first-stop paths from a separate diagonal sequence tend to
release (square).
(c) The fixed local patch crossed strictly at
$X_*=\mathcal C_{\lambda,s}(\alpha_0)$ selects that prefix along the
same sequence as (a). The dotted circle indicates localization only;
the physical patch and its separation from earlier prefixes are
established in full state space. Other coordinates and distinct
incoming arcs are suppressed, not identified.}
\label{fig:observation-rules}
\end{figure}


\subsection{Common paths beyond a complete contact}

Write $f_\eps^\tau$ for the actual path stopped at its first crossing
of $p=-a\tau$, with the same normalization of its parameter interval
as in \eqref{eq:ps5-stop}. The return curve and phase are unchanged. Since $F_\lambda(s)>s$,
there is $s<\sigma\le s_0$ with $\sigma<F_\lambda(s)$ for every
$0<s<s_0$. After a common restriction with $2s\le s_0$ and
$F_\lambda(s)>2s$, one may take $\sigma=2s$.

\begin{corollary}[Common paths beyond a complete contact]
\label{cor:common-downstream-contact}
Fix $\lambda$ and $0<s<\sigma\le s_0$ with
$\sigma<F_\lambda(s)$. The common sequence \eqref{eq:ps5-common}
with target $d=s$ satisfies
\[
 \sup_{(k,\theta,\nu,\mathbf n)}
 \dF\bigl(f_{\eps_m(s)}^\sigma,
  \cW_{\lambda,\nu,\sigma}(s,j_\sigma(s))\bigr)\longrightarrow0.
\]
The corresponding image convergence is uniform in $\dH$.
Here $2<j_\sigma(s)<3$, so the common limiting path includes
the entire contact $\cC_{\lambda,s}|_{[0,2]}$ and a positive-length
prefix of its following free flight.
\end{corollary}

\begin{proof}
The inequalities on $\sigma$ give
$F_\lambda^{-1}(\sigma)<s<\sigma$. Thus $d=s$ is an interior
free target for the $\sigma$ section. Apply
Theorem~\ref{thm:common-free} at that section. Its sequence depends
on $\lambda,d$, but not on the section, so it is precisely
\eqref{eq:ps5-common} with $d=s$. Strict decrease of $p$ along
the free flight and $s<\sigma<F_\lambda(s)$ give
$2<j_\sigma(s)<3$. The flight duration $T_s$ is positive, and
the retained free duration $(j_\sigma(s)-2)T_s$ is therefore positive.
\end{proof}

\subsection{First-angle truncation}

For an oriented path $g:[0,1]\to\R^5$ that starts above and reaches
the section $p=-as$, define its first-hit truncation by
\[
 \mathcal T_s g
   =g|_{[0,\tau_s(g)]},\qquad
 \tau_s(g)=\min\{q\in[0,1]:p(g(q))=-as\},
\]
reparameterized onto $[0,1]$. This operation is well defined in the
zero-distance quotient of \eqref{eq:ps5-frechet}. Indeed, two
zero-distance paths admit a common representation
$g\circ a=\widetilde g\circ b$, with $a,b$ continuous
nondecreasing surjections. Such a representation is obtained by
taking a limit of the increasing comparison graphs and
parameterizing the limiting monotone relation by the sum of its
coordinates. At the first section
point of the common representation, $a$ and $b$ reach the first
section parameters of $g$ and $\widetilde g$, respectively.
Restricting $a,b$ there gives a common representation of their
prefixes. Thus only inserted pauses can differ, and the prefixes
have zero ordered distance.

\begin{corollary}[Discontinuity of first-hit truncation]
\label{cor:first-hit-discontinuity}
Under the same restrictions on $\lambda,s,\sigma$, fix a complete
label $\mathfrak l=(\lambda,k,\theta,\nu,\mathbf n)$ and put
\[
 \Gamma=\cW_{\lambda,\nu,\sigma}(s,j_\sigma(s)).
\]
There is a sequence $\eps_\ell\downarrow0$ for this label such that
\[
 \begin{split}
 \dF(f_{\eps_\ell}^\sigma,\Gamma)&\longrightarrow0,\\
 \dF(\mathcal T_s f_{\eps_\ell}^\sigma,
       \cW_{\lambda,\nu,s}(s,2))&\longrightarrow0,
 \end{split}
\]
whereas
$\mathcal T_s\Gamma=\cW_{\lambda,\nu,s}(s,0)$ in the ordered
quotient. These two stopped limit paths have different endpoints.
Consequently $\mathcal T_s$ is discontinuous at $\Gamma$, even on
the subspace consisting of $\Gamma$ and actual paths stopped at
the $\sigma$ section.
\end{corollary}

\begin{proof}
Choose $d_\ell\uparrow s$ with
$F_\lambda^{-1}(\sigma)<d_\ell<s$. Each $d_\ell$ is an interior
free target for both the $s$ and $\sigma$ sections. The explicit
sequences supplied by Theorem~\ref{thm:common-free} at these two
sections are identical:
\[
 \eps_m(d_\ell)=
 \exp\{-\tfrac{32}{9}(3/4)^{-m}\mathcal S_\lambda(d_\ell)\}.
\]
For each $\ell$, take $m_\ell$ sufficiently large that both
ordered-path errors are below $1/\ell$, and also that
$\eps_{m_\ell}(d_\ell)<1/\ell$ and, when $\ell>1$,
$\eps_{m_\ell}(d_\ell)<\eps_{m_{\ell-1}}(d_{\ell-1})/2$.
Write $\eps_\ell=\eps_{m_\ell}(d_\ell)$.

Continuous dependence of the finite-scale return curves and their
free-section roots gives
\[
 j_\sigma(d_\ell)\longrightarrow j_\sigma(s),\qquad
 j_s(d_\ell)\longrightarrow j_s(s)=2.
\]
The complete-path continuity proved after
\eqref{path:whole-error}, applied separately at the fixed sections
$s$ and $\sigma$, therefore gives
\[
 \begin{split}
 \cW_{\lambda,\nu,\sigma}(d_\ell,j_\sigma(d_\ell))
      &\longrightarrow\Gamma,\\
 \cW_{\lambda,\nu,s}(d_\ell,j_s(d_\ell))
      &\longrightarrow\cW_{\lambda,\nu,s}(s,2).
 \end{split}
\]
The two diagonal error bounds prove the two asserted actual limits.
The actual $s$ section is reached before the $\sigma$ section,
and its prefix is exactly
$\mathcal T_s f_{\eps_\ell}^\sigma=f_{\eps_\ell}^s$ in the
ordered quotient.

The incoming arc has $p\ge0$. Every return preceding the base $s$
in $L_s$ has $p>-as$ except that the last reaches $-as$ at its
final touchdown, the admission point $\cC_{\lambda,s}(0)$.
Hence $\Gamma$ first reaches this section at that admission point,
and its first-hit prefix is $\cW_{\lambda,\nu,s}(s,0)$.
The admission slip is $U_s>0$ from \eqref{gr:embedding}, while
the release slip is zero. In particular,
\[
 |\cC_{\lambda,s}(2)-\cC_{\lambda,s}(0)|
   \ge \frac{U_s}{|\cos\phi_s|}>0.
\]
The ordered distance bounds the distance between terminal states,
so the limits of the truncated paths and of the long paths cannot
commute. This proves the claimed discontinuity.
\end{proof}

\subsection{Fixed local observation patches}

The angular section contains the whole limiting contact. A local patch
can cross it at any prescribed interior point: the next construction
uses slip on compression, normal velocity on release, and the two
one-sided crossings at arrest.

\begin{corollary}[Common contact prefixes at a local observation surface]
\label{cor:common-local-contact}
Fix $\lambda$, $0<s<\sigma\le s_0$ with
$\sigma<F_\lambda(s)$, and $\alpha_0\in(0,2)$. Put
\[
 X_*=\cC_{\lambda,s}(\alpha_0),\qquad
 h_s=\mu A(\phi_s)-B(\phi_s)>0,
\]
and choose $c_\Sigma>as/h_s$. Define on physical states
\[
 u_0(\cX)=\omega\cos\phi+U_*,\qquad
 \Psi(\cX)=v-c_\Sigma u_0(\cX).
\]
There is $r_\Sigma>0$ such that the closed observation patch
\[
 \Sigma=\{\cX:|\cX-X_*|\le r_\Sigma,\quad
                         \Psi(\cX)=\Psi(X_*)\}
\]
has the following property. For all sufficiently large $m$, every
actual trajectory at $\eps=\eps_m(s)$ in the prescribed compact auxiliary label class
first meets $\Sigma$ before reaching $p=-a\sigma$, and this first
meeting occurs during the actual contact whose base tends to $s$.
If $f_\eps^\Sigma$ is the path stopped at this first meeting and
reparameterized onto $[0,1]$, then
along the common sequence \eqref{eq:ps5-common} with $d=s$,
\[
 \sup_{(k,\theta,\nu,\mathbf n)}
 \dF\bigl(f_{\eps_m(s)}^\Sigma,
             \cW_{\lambda,\nu,s}(s,\alpha_0)\bigr)
       \longrightarrow0.
\]
The image convergence is also uniform in $\dH$. The patch is
independent of $\eps$ and the auxiliary labels; each incoming datum
retains its own arc $I_{\lambda,\nu}$.
\end{corollary}

\begin{proof}
The exact contact formulas \eqref{gr:embedding} and
\eqref{eq:ps2-reference-path} give, on compression,
\[
 \frac{\dd}{\dd\alpha}\Psi(\cC_{\lambda,s}(\alpha))
       =(c_\Sigma-as/h_s)U_s>0\qquad(0<\alpha<1).
\]
On release $u_0=0$, and \eqref{loc:normal-reference} gives
\[
 \frac{\dd}{\dd\alpha}\Psi(\cC_{\lambda,s}(\alpha))
   =\tfrac12s^{3/2}R_s
      (Z_s-Y_s)((\alpha-1)R_s)>0\qquad(1<\alpha<2).
\]
The active force is positive before its first zero. These two
formulas and the continuous join show that $\Psi$ is strictly
increasing along the entire reference contact. At $\alpha_0=1$
we use the two one-sided strict crossings.

The compact tail $L_s$ meets $p=-as$ only at
$\cC_{\lambda,s}(0)$, whose $\Psi$ value is smaller than
$\Psi(X_*)$. Hence $X_*\notin L_s$. Choose $r_\Sigma>0$ so that the
closed $2r_\Sigma$ ball about $X_*$ avoids $L_s$ and is contained in
$\{p<0\}$. Every incoming arc lies in $\{p\ge0\}$, so this
choice also separates the patch from all incoming arcs uniformly.

Choose $\eta>0$ with
$0<\alpha_-:=\alpha_0-\eta<\alpha_+:=\alpha_0+\eta<2$
so small that $\cC_{\lambda,s}([\alpha_-,\alpha_+])$
lies in the open $r_\Sigma$ ball. At the two ends $\Psi-\Psi(X_*)$
has opposite signs. The earlier contact segment
$\cC_{\lambda,s}([0,\alpha_-])$ has strictly smaller
$\Psi$ value and is a positive distance from the closed patch.
Together with the preceding tail and incoming-arc separation,
this gives a uniform positive distance between $\Sigma$ and
the entire preceding reference prefix.

Corollary~\ref{cor:common-downstream-contact} and its ordered
pairings give uniform convergence of the actual paths through
this contact. More specifically,
Lemma~\ref{path:ordered-return} pairs its compression and release
with $\cC_{\lambda,s}|_{[0,2]}$, after using the convergence of
the actual base to $s$. These pairings align the earlier prefixes
as in the proof of Theorem~\ref{thm:whole}. Thus, for large $m$,
the actual trajectory cannot meet $\Sigma$ before the paired
point $\alpha_-$, while the whole paired bracket stays in the
open $r_\Sigma$ ball and its end values of $\Psi-\Psi(X_*)$ have opposite
signs. The intermediate value theorem supplies a meeting inside
this bracket. Since $\Sigma$ is closed, the nonempty set of
meeting times up to the end of the bracket is compact and has a
first element. It lies between the paired points $\alpha_-$
and $\alpha_+$, within this actual contact.

Restrict the ordered pairings at that first meeting. Their
errors tend to zero uniformly in the auxiliary labels, and the
reference stopping parameter lies in
$[\alpha_0-\eta,\alpha_0+\eta]$. Uniform continuity of
$\cC_{\lambda,s}$, followed by $\eta\downarrow0$, proves the
asserted prefix convergence. The image assertion follows from
$\dH\le\dF$.
\end{proof}

Appendix~\ref{sec:accumulation-geometry} identifies the single
terminal-parameter duplication and proves that the full-state accumulation
families are circles. Appendix~\ref{sec:memory} records the interval of
scaled effective-gap velocities above the release endpoint path.
